\documentclass{situationist-preprint}

\title{Disagreement and Moral Realism}
\msid{MTE-12-a}
\msfamily{MTE-12}
\msdiscipline{Metaethics}
\mstype{Principal manuscript}
\mskeywords{moral disagreement; argument from relativity; moral realism; scientific disagreement; idealization; peer disagreement}

\begin{document}
\maketitle

\begin{abstract}
The argument from disagreement holds that the extent and persistence of moral disagreement is better explained by the absence of moral facts than by the hypothesis that some parties are mistaken. We argue that moral disagreement is no deeper than disagreement in other theoretical domains, such as physics, history and mathematics, that are not thought to lack facts. We distinguish three measures of the depth of a disagreement: its persistence among informed and competent inquirers, its resistance to explanation by distorting factors, and its location at the level of fundamental principle rather than application. On each measure, moral disagreement is matched by disagreement in domains whose factuality nobody doubts. Most moral disagreement is explained by disagreement about non-moral facts, by partiality, and by inconsistent application of shared principles; the residue that survives idealization is comparable in extent to the residue in the foundations of physics or in the interpretation of history. An argument that inferred the absence of facts from disagreement of this depth would prove too much. Disagreement therefore gives no evidence against realism.
\end{abstract}

\section{The problem}

\citet[ch.~1, \S 8]{Mackie1977} argued that the variation in moral codes between societies and periods is more readily explained by the hypothesis that moral codes reflect ways of life than by the hypothesis that they express perceptions, most of them seriously inadequate and badly distorted, of objective values. The argument from relativity, or from disagreement, is an inference to the best explanation: the hypothesis that there are no moral facts explains the data at least as well as realism does, and more economically.

The argument has been developed in two directions. The first is empirical: \citet{DorisPlakias2008} argue that some moral disagreements, such as those between members of cultures of honor and others about the permissibility of violent responses to insult, would persist under idealized conditions, so that the realist cannot attribute them to ignorance or irrationality. The second is epistemological: if equally competent inquirers disagree, conciliatory norms of rational belief require each to suspend judgment, and moral knowledge becomes unattainable even if there are moral facts. This version is not an argument against realism but against moral knowledge, and the conclusion of EPI-07, that one should conciliate with an epistemic peer while giving some additional weight to one's own first-order evidence, implies that its force depends on how much moral disagreement is genuinely peer disagreement. We return to it in Section~\ref{s:obj}.

Realists have replied in two ways. \citet{Brink1984} argued that much moral disagreement is explained by disagreement about non-moral facts, by the distorting influence of self-interest and ideology, and by failures of imagination, and that the residue is no greater than realism can accommodate. \citet{Enoch2009} asked how, exactly, disagreement is supposed to be a problem for realism, and argued that none of the available ways of turning it into an argument succeeds without premises that would also condemn other domains. We develop the second reply, by specifying what it would be for a disagreement to be deep and showing that moral disagreement is not deep in any sense in which disagreement in other theoretical domains is not.

\section{Measures of depth}

\begin{definition}[Depth]\label{d:depth}
A disagreement about a question is \emph{deep} to the extent that:
\begin{enumerate}[label=(\alph*)]
\item \emph{persistence}: it persists among inquirers who share the relevant evidence, are competent in the methods of the domain, and have engaged with each other's arguments;
\item \emph{inexplicability}: it is not explained by distorting factors, such as ignorance of relevant facts, partiality, inattention, errors of reasoning or inconsistent application of principles the parties share;
\item \emph{fundamentality}: it concerns the basic principles of the domain rather than their application to particular cases.
\end{enumerate}
\end{definition}

The argument from disagreement requires that moral disagreement be deep in all three respects. A disagreement that disappears under shared evidence is evidence of ignorance. One explained by distorting factors is evidence of those factors. One confined to applications is consistent with agreement on principles, and so with the existence of the facts that the principles state.

\begin{principle}[Parity]\label{p:parity}
If a disagreement in one domain is no deeper, on each measure of Definition~\ref{d:depth}, than a disagreement in another domain, then the first is no stronger evidence that its domain lacks facts than the second is that its domain lacks facts.
\end{principle}

Principle~\ref{p:parity} states the evidential symmetry on which the argument from disagreement itself relies. The argument infers the absence of facts from features of disagreement; if those features are present to the same degree elsewhere, the inference has the same strength elsewhere.

\section{The comparison}

\paragraph{Persistence.} Physicists who share all the evidence and have engaged with each other's arguments disagree about the interpretation of quantum mechanics: whether the wave function collapses, whether there are many worlds, whether hidden variables determine outcomes. The disagreement has persisted for nearly a century among the most competent inquirers in the field. Historians disagree about the causes of the First World War and of the fall of the Western Roman Empire, with access to the same archives. Mathematicians disagree about whether the continuum hypothesis has a determinate truth value and about which new axioms of set theory to accept. Economists disagree about the effects of minimum wages and of monetary policy. None of these disagreements is taken to show that there is no fact about the wave function, the causes of a war, or the effects of a policy. The persistence of moral disagreement among competent moral philosophers about, say, the moral status of the fetus or the permissibility of consequentialist trade-offs is of the same order.

\paragraph{Inexplicability.} Most moral disagreement is explained by distorting factors. Disagreement about capital punishment is in large part disagreement about deterrence, about the reliability of convictions and about the distribution of error, which are non-moral questions. Disagreement about the treatment of animals has been in large part disagreement about their capacities, as the evidence on which ETH-07 relies makes clear. The historical disagreements that Mackie cites, between slaveholding and non-slaveholding societies, were sustained by false beliefs about the capacities of the enslaved and by the self-interest of those who benefited, and they did not survive the correction of either. Disagreements across cultures often reflect the same principle applied to different circumstances: societies that permit practices we condemn frequently hold false beliefs about their effects, or face conditions of scarcity under which the practices serve ends we share. Doris and Plakias's culture of honor is no exception. The norms of honor they describe are explained by the conditions of herding economies without effective policing, in which a reputation for retaliation protects property; those who hold the norms share with others the principle that one may protect oneself and one's dependants, and differ in their beliefs about what protection requires. In each case, the disagreement is explained by factors that idealization removes: information, impartiality and consistency.

Disagreement in other domains is explained by distorting factors to a similar extent. Much scientific disagreement is sustained by differences in the evidence the parties have attended to, by loyalty to research programs, and by the interests of those who fund research. The residue that is not so explained, such as the disagreement over the interpretation of quantum mechanics, is the analogue of the residue in ethics.

\paragraph{Fundamentality.} Moral disagreement is, for the most part, disagreement about application. Every moral tradition of which we have record condemns gratuitous cruelty, regards some promises as binding, gives weight to the welfare of children, and distinguishes the innocent from the guilty. Disagreement concerns the scope of these principles, which beings fall under them, and how they are to be weighed against each other in hard cases. Disagreement about fundamentals, of the kind that divides consequentialists and deontologists, is confined to theorists, and ETH-02 argues that even it is less fundamental than it appears, since every deontological theory is extensionally equivalent to a consequentialist theory with a suitable value function. Physics, by contrast, has disagreement about fundamentals at its frontier: whether spacetime is fundamental, whether physical law is deterministic. Moral disagreement is, if anything, less fundamental than disagreement in physics.

\begin{theorem}\label{t:main}
Moral disagreement is no deeper than disagreement in other theoretical domains that are not thought to lack facts, and so, by Principle~\ref{p:parity}, it is no evidence that morality lacks facts.
\end{theorem}

\begin{proof}
By the comparison of this section, on each measure of Definition~\ref{d:depth} there is a domain not thought to lack facts, physics, history, mathematics or economics, in which disagreement is at least as deep as moral disagreement. If moral disagreement were evidence that morality lacks facts, then by Principle~\ref{p:parity} disagreement in those domains would be equal evidence that they lack facts. It is not, since we have overwhelming reason to believe that there are facts about the causes of wars and the behavior of particles. So moral disagreement is no such evidence.
\end{proof}

Where moral disagreement would persist under idealization, the right conclusion is the one drawn in other domains: the question may lack the determinate answer the parties seek, or the evidence may underdetermine it. MTE-01 holds that moral facts are natural facts about idealized reactions, and that where idealized reactions diverge neither party's verdict that an act is forbidden or required is true, and the comparative question between the options has no answer. That is a local indeterminacy, of the kind that physics tolerates in the boundaries of its kinds and history in the individuation of its events, and it is no wider in ethics than elsewhere.

\section{Objections and replies}\label{s:obj}

\begin{objection}[Convergence in science]
Scientific disagreements are eventually resolved, and science exhibits convergence over time. Moral disagreements are not resolved in the same way.
\end{objection}

\begin{reply}
Moral disagreements are resolved: about slavery, about the subordination of women, about the torture of suspects as a routine part of criminal procedure, about the burning of heretics. MTE-19 argues that moral progress is real and measurable as the expansion of the moral circle, though neither steady nor irreversible. Scientific convergence is also neither steady nor irreversible, and the frontier of every science is a region of disagreement. The objection compares the settled core of science with the frontier of ethics.
\end{reply}

\begin{objection}[Idealized persistence]
Some moral disagreements would persist among fully informed, impartial and consistent agents \citep{DorisPlakias2008}. No comparable claim holds for physics, where idealized inquirers with all the evidence would agree.
\end{objection}

\begin{reply}
It is not known that idealized physicists would agree about the interpretation of quantum mechanics, since the interpretations are empirically equivalent on all the evidence we have and may be on all the evidence there could be. If there are disagreements that would persist under idealization in ethics, there are disagreements that would persist under idealization in physics and in mathematics, where the independence of the continuum hypothesis from the standard axioms allows idealized set theorists to differ. The response in each domain is the same: the question is underdetermined, or there is no fact, locally. Nobody infers from the case of the continuum hypothesis that there are no facts about the natural numbers.
\end{reply}

\begin{objection}[Sceptical version]
Even if disagreement is no evidence against moral facts, it is evidence against moral knowledge: given disagreement with peers, conciliation requires suspension of judgment.
\end{objection}

\begin{reply}
The sceptical version, if sound, would also yield scepticism about the interpretation of quantum mechanics and the causes of wars, which is the right result for contested questions in every domain and leaves uncontested moral knowledge untouched. Since most moral disagreement is explained by distorting factors, most of it is not peer disagreement: one party lacks information, imagines less fully, or applies its principles inconsistently. On the view of EPI-07, conciliation is required only toward peers, and even then leaves additional weight for one's own first-order evidence. Disagreement therefore leaves intact our knowledge that gratuitous cruelty is wrong, and lowers our confidence in contested verdicts, as it should \citep{McGrath2008,Vavova2014}.
\end{reply}

\section{Conclusion}

The argument from disagreement requires that moral disagreement be deep in a way that disagreement in factual domains is not. It is not. Moral disagreement persists among competent inquirers no more than disagreement in physics, history and mathematics; it is explained by distorting factors at least as often; and it concerns application more than fundamental principle. Since the argument would show, with equal force, that there are no facts about the causes of wars or the behavior of particles, moral disagreement is no evidence against moral realism.

\begin{thebibliography}{99}\small
\bibitem[Brink(1984)]{Brink1984} Brink, D.~O. (1984). Moral realism and the sceptical arguments from disagreement and queerness. \emph{Australasian Journal of Philosophy}, 62(2), 111--125.
\bibitem[Doris and Plakias(2008)]{DorisPlakias2008} Doris, J.~M. and Plakias, A. (2008). How to argue about disagreement: Evaluative diversity and moral realism. In W. Sinnott-Armstrong (ed.), \emph{Moral Psychology, Volume 2: The Cognitive Science of Morality}. MIT Press.
\bibitem[Enoch(2009)]{Enoch2009} Enoch, D. (2009). How is moral disagreement a problem for realism? \emph{The Journal of Ethics}, 13(1), 15--50.
\bibitem[Mackie(1977)]{Mackie1977} Mackie, J.~L. (1977). \emph{Ethics: Inventing Right and Wrong}. Penguin.
\bibitem[McGrath(2008)]{McGrath2008} McGrath, S. (2008). Moral disagreement and moral expertise. \emph{Oxford Studies in Metaethics}, 3.
\bibitem[Vavova(2014)]{Vavova2014} Vavova, K. (2014). Moral disagreement and moral skepticism. \emph{Philosophical Perspectives}, 28.
\end{thebibliography}

\end{document}
